\section{Bruk av kunstig intelligens}
\label{sec:ai-bruk}

KI-verktøy, primært chat-baserte språkmodeller, ble brukt
aktivt gjennom prosjektet. Siden ingen på teamet hadde
erfaring med C\# eller .NET fra før, fungerte KI som et
støtteverktøy for å forstå rammeverkskonsepter som dependency
injection, \texttt{DelegatingHandler}-mønsteret og multi-stage
Dockerfiler.

I implementasjonsarbeidet ble KI brukt til generering av
rutinekode, diskusjon av eksisterende kode, feilsøking og
generering av testscenarier ut fra teamets spesifikasjoner.
Det mest omfattende eksempelet er ende-til-ende-testskriptet
omtalt i delkapittel~\ref{sec:e2e-testing}, som ble manuelt
validert mot reelle testmiljøer før det ble brukt.

KI-forslag ble ikke akseptert ukritisk. Særlig ved integrasjon
mot AntMedia og NHN var domenekunnskap om faktisk API-oppførsel
avgjørende, og forslag som ikke samsvarte med dokumentasjon,
kravspesifikasjon eller testet oppførsel ble forkastet.
Arkitektur, designvalg og kjernelogikk er resultatet av teamets
egne vurderinger.

I rapportarbeidet ble KI brukt til \LaTeX-hjelp, språkkorrektur
og forslag til struktur i tekst gruppen selv hadde skrevet.
KI er ikke brukt som faktakilde. Alle faglige påstander er
forankret i siterte kilder, kravspesifikasjonen eller egen
empirisk testing. Bruken er også redegjort for i
KI-deklarasjonen i vedlegg~\ref{app:ki-deklarasjon}.